\section{Analysis Metric 5: Unique Ethereum Addresses}
\label{sec:analysisEthAdd}

\subsection{Metric Definition}
\label{sec:metricEthAddDefinition}

The last metric we analyzed is the number of unique Ethereum addresses
found in the plain text files of the \ac{NPM} packages.

To compute this metric and identify strings that represent Ethereum addresses, 
the developed tool (\cref{sec:ToolForAnalyzingNpmPackages}) uses the following precompiled regular expression:
\begin{verbatim}
(?<![\w\-/_])0x[a-fA-F0-9]{40}(?![\w\-/_])
\end{verbatim}

The regex uses a lookbehind \texttt{(?<![{\textbackslash}w{\textbackslash}-/\_])}
and a lookahead \texttt{(?![{\textbackslash}w{\textbackslash}-/\_])}
to ensure that the address is neither preceded nor followed by words,
hyphens, underscores or slashes.
This avoids, for example, capturing a string that is part of a URL or a file path,
reducing false positives.
Unlike word boundaries \texttt{{\textbackslash}b},
this solution also correctly handles the typical path separators, like slashes and hyphens.

A valid Ethereum address consists of exactly 40 hexadecimal characters
(i.e., digits 0 to 9 and letters a to f, both lowercase and uppercase) 
preceded by the prefix \texttt{0x}~\cite{ethereum}.

The tool identifies all occurrences of this pattern in each plain text file of a package version,
ignores duplicates within that version to ensure the values are unique, and then 
sums the counts to obtain the total number of unique Ethereum addresses for the entire package version.

This metric was taken into consideration because attacks B (\cref{sec:attackB}) 
introduce seven smart contract Ethereum addresses.


\subsection{Goodware Dataset Scenario}
\label{sec:goodwareEthAdd}

For this metric, the Goodware dataset (\cref{sec:goodwareDataset}) contains the 
total count of unique Ethereum addresses for each of the 20 versions of every package.

For this metric we do not present any figure because in the Goodware dataset only
two packages have exactly one Ethereum address in at least one version, namely
\texttt{brace-expansion} and \texttt{dotenv-expand}.
We analyze both packages to understand the reason for this presence.

\begin{itemize}
    \item \textbf{brace-expansion:} This package provides brace expansion for JavaScript,
    allowing the use of shell-like syntax to generate sets of strings from a pattern~\cite{brace-expansion}.
    
    Starting from version 4.0.0 (the second-to-last available version), the tarball includes
    the file \texttt{tea.yaml}, a configuration file required to register an
    open-source project on the Tea protocol.
    The Tea protocol is an incentivization protocol for open-source projects,
    which allows registered projects to earn TEA tokens based on their popularity and
    usage~\cite{teaProtocol}.
    The \texttt{tea.yaml} file must contain one or a list of Ethereum addresses
    to receive TEA tokens and to define who controls the project.
    The file is also present in the latest version 4.0.1.

    \item \textbf{dotenv-expand:} This package extends the functionality of dotenv by 
    allowing variable expansion in dotenv files~\cite{dotenv-expand}.

    Similarly to \texttt{brace-expansion}, 
    starting from version 11.0.7 the tarball includes the file \texttt{tea.yaml}
    containing an Ethereum address.
    This file remained present for four consecutive versions, until 
    version 12.0.3, the most recent available, in which it was removed.

\end{itemize}

\subsection{Comparison with Malware Dataset}
\label{sec:analysisEthAddComparisonMalware}

Although the Goodware dataset (\cref{sec:goodwareDataset}) contains very few packages with Ethereum addresses,
as discussed in \cref{sec:goodwareEthAdd}, 
the metric nonetheless proves effective in detecting attack B (\cref{sec:attackB}).
The compromised versions of this attack introduce seven known smart contract Ethereum addresses into the packages, 
which the tool (\cref{sec:ToolForAnalyzingNpmPackages}) correctly identifies.
Given that only two packages in the Goodware dataset contain even a single Ethereum address, 
a count of seven qualifies the malicious versions as outliers.

However, the metric is not indicative for the remaining three attacks.
For attack A (\cref{sec:attackA}), the malicious \ac{DLL} is a binary file 
and is therefore excluded from this computation of the metric, 
leaving the Ethereum addresses count unchanged with respect to the legitimate versions.
For attacks C (\cref{sec:attackC}) and D (\cref{sec:attackD}),
the malicious code targets the theft of credentials and cloud services access keys, 
rather than manipulating cryptocurrency transactions.
As a result, neither \texttt{bundle.js} nor \texttt{bun\_environment.js} 
contains any Ethereum addresses, and the count remains unchanged compared to the legitimate versions.
In all three cases, the compromised versions are therefore indistinguishable from the legitimate ones under this metric.